% Physical page 2 of the concise study, which carries this sheet and nothing
% else. The dates, relation labels, alternatives and unresolved states are the
% generated projection of research/chronology.toml, and stand exactly as the
% reviewed appendix of the expansive study prints them; the explanatory rows
% condense that appendix's. The generated definitions are imported here, once,
% because the concise entrypoint's own preamble does not carry them.
% Generated. Do not edit.
% Deterministic projection of research/chronology.toml.
% schema: 4
% document: liturgy/roman-rite/1962/propers/temporal/60-twentieth-after-pentecost
% generated-by: tools/tpt proper-chronology annotations
\newcommand{\chronologyannotationclaim}[6]{#6}
\newcommand{\chronologyannotationcomparisonclaim}[8]{#8}
\newcommand{\chronologyannotationanchorclaim}[9]{#9}
\newcommand{\chronologyannotationanchorgroup}[8]{#8}
\newcommand{\chronologyannotationreach}[3]{}
\newcommand{\chronologyannotationgroup}[3]{#3}
\newcommand{\chronologyannotationcomparisongroup}[4]{#4}
\newcommand{\chronologyannotation}[1]{%
  \ifcsname triptychchronologyannotation@#1\endcsname
    \csname triptychchronologyannotation@#1\endcsname
  \else
    \PackageError{triptych}{No chronology annotation for #1}{}%
  \fi}
% comparison-dependency: src/claude/liturgy/roman-rite/1962/propers/temporal/60-twentieth-after-pentecost/research/chronology-profile-comparisons.toml sha256=139a40c7f50531ff34337f64368f66dc23e9909ef9d06f2802237a8200eee01c
% introit: Introit
\expandafter\def\csname triptychchronologyannotation@introit\endcsname{%
  \chronologyannotationgroup{historical-setting}{nonuniform}{\textbf{Historical setting} -- No single occasion date applies across every cited locus.}%
  \space \chronologyannotationgroup{composition}{nonuniform}{\textbf{Composition} -- No single date applies across every cited locus.}%
}
% epistle: Epistle
\expandafter\def\csname triptychchronologyannotation@epistle\endcsname{%
  \chronologyannotationgroup{composition}{disputed}{\textbf{Composition} -- disputed: \chronologyannotationreach{Eph.5.15}{Eph}{true}\chronologyannotationreach{Eph.5.16}{Eph}{true}\chronologyannotationreach{Eph.5.17}{Eph}{true}\chronologyannotationreach{Eph.5.18}{Eph}{true}\chronologyannotationreach{Eph.5.19}{Eph}{true}\chronologyannotationreach{Eph.5.20}{Eph}{true}\chronologyannotationreach{Eph.5.21}{Eph}{true}\chronologyannotationclaim{composition.epistle-to-the-ephesians}{composition}{catholic-traditional-v1}{disputed}{a period between 58 and 63}{A.D. 58--63}; \chronologyannotationreach{Eph.5.15}{Eph}{true}\chronologyannotationreach{Eph.5.16}{Eph}{true}\chronologyannotationreach{Eph.5.17}{Eph}{true}\chronologyannotationreach{Eph.5.18}{Eph}{true}\chronologyannotationreach{Eph.5.19}{Eph}{true}\chronologyannotationreach{Eph.5.20}{Eph}{true}\chronologyannotationreach{Eph.5.21}{Eph}{true}\chronologyannotationclaim{composition.epistle-to-the-ephesians}{composition}{catholic-traditional-v1}{disputed}{(Philemon; Colossians; Ephesians; Philippians), 61}{A.D. 61}.}%
  \space \chronologyannotationcomparisongroup{catholic-critical-v1}{composition}{composition-only}{\textbf{Composition (Catholic critical chronology, for comparison)} -- disputed: \chronologyannotationreach{Eph.5.15}{Eph}{true}\chronologyannotationreach{Eph.5.16}{Eph}{true}\chronologyannotationreach{Eph.5.17}{Eph}{true}\chronologyannotationreach{Eph.5.18}{Eph}{true}\chronologyannotationreach{Eph.5.19}{Eph}{true}\chronologyannotationreach{Eph.5.20}{Eph}{true}\chronologyannotationreach{Eph.5.21}{Eph}{true}\chronologyannotationcomparisonclaim{critical.ephesians-later-disciple}{composition}{catholic-critical-v1}{catholic-critical-v1}{disputed}{passage.united-states-conference-of-catholic-bishops.new-american-bible-revised-edition.english-usccb-web-2026-09-28.ephesians-introduction}{around A.D. 80--100}{Ephesians under the later-disciple hypothesis (approximate), c. A.D. 80--100}.}%
}
% gradual: Gradual
\expandafter\def\csname triptychchronologyannotation@gradual\endcsname{%
  \chronologyannotationgroup{traditional-attribution}{disputed}{\textbf{Traditional attribution} -- disputed: \chronologyannotationreach{Ps.144.15}{Ps.107 Ps.144}{true}\chronologyannotationreach{Ps.144.16}{Ps.107 Ps.144}{true}\chronologyannotationclaim{israel.monarchy.david-traditional-era}{traditional-attribution}{catholic-traditional-v1}{disputed}{reigned from 1055 to 1015 B.C.}{David (reign in the usual chronology), B.C. 1055--1015}.}%
  \space \chronologyannotationgroup{historical-setting}{research-pending}{\textbf{Historical setting} -- Occasion date unresolved.}%
  \space \chronologyannotationgroup{composition}{preferred}{\textbf{Composition}: \chronologyannotationreach{Ps.144.15}{Ps}{true}\chronologyannotationreach{Ps.144.16}{Ps}{true}\chronologyannotationclaim{critical.psalms.latest-composition-boundary}{composition}{catholic-critical-v1}{preferred}{before the Maccabean period, around 165 B.C.; no individual psalm can be dated securely}{Before c. 165 B.C.}}%
}
% alleluia: Alleluia
\expandafter\def\csname triptychchronologyannotation@alleluia\endcsname{%
  \chronologyannotationgroup{traditional-attribution}{disputed}{\textbf{Traditional attribution} -- disputed: \chronologyannotationreach{Ps.107.2}{Ps.107 Ps.144}{true}\chronologyannotationclaim{israel.monarchy.david-traditional-era}{traditional-attribution}{catholic-traditional-v1}{disputed}{reigned from 1055 to 1015 B.C.}{David (reign in the usual chronology), B.C. 1055--1015}.}%
  \space \chronologyannotationgroup{historical-setting}{research-pending}{\textbf{Historical setting} -- Occasion date unresolved.}%
  \space \chronologyannotationgroup{composition}{preferred}{\textbf{Composition}: \chronologyannotationreach{Ps.107.2}{Ps}{true}\chronologyannotationclaim{critical.psalms.latest-composition-boundary}{composition}{catholic-critical-v1}{preferred}{before the Maccabean period, around 165 B.C.; no individual psalm can be dated securely}{Before c. 165 B.C.}}%
}
% gospel: Gospel
\expandafter\def\csname triptychchronologyannotation@gospel\endcsname{%
  \chronologyannotationgroup{narrated-event}{preferred}{\textbf{Event}: \chronologyannotationreach{John.4.46}{John.4.46-53}{false}\chronologyannotationreach{John.4.47}{John.4.46-53}{false}\chronologyannotationreach{John.4.48}{John.4.46-53}{false}\chronologyannotationreach{John.4.49}{John.4.46-53}{false}\chronologyannotationreach{John.4.50}{John.4.46-53}{false}\chronologyannotationreach{John.4.51}{John.4.46-53}{false}\chronologyannotationreach{John.4.52}{John.4.46-53}{false}\chronologyannotationreach{John.4.53}{John.4.46-53}{false}\chronologyannotationclaim{life-of-christ.healing-of-the-rulers-son-at-cana}{narrated-event}{catholic-traditional-v1}{preferred}{Passover, A.U.C. 779 - about Pentecost, 780}{The healing of the ruler's son at Cana, A.D. 26--27 (derived)}.}%
  \space \chronologyannotationgroup{composition}{disputed}{\textbf{Composition} -- disputed: \chronologyannotationreach{John.4.46}{John}{true}\chronologyannotationreach{John.4.47}{John}{true}\chronologyannotationreach{John.4.48}{John}{true}\chronologyannotationreach{John.4.49}{John}{true}\chronologyannotationreach{John.4.50}{John}{true}\chronologyannotationreach{John.4.51}{John}{true}\chronologyannotationreach{John.4.52}{John}{true}\chronologyannotationreach{John.4.53}{John}{true}\chronologyannotationclaim{composition.gospel-of-john}{composition}{catholic-traditional-v1}{disputed}{from the year 90 to 100 (approximately)}{c. A.D. 90--100}; \chronologyannotationreach{John.4.46}{John}{true}\chronologyannotationreach{John.4.47}{John}{true}\chronologyannotationreach{John.4.48}{John}{true}\chronologyannotationreach{John.4.49}{John}{true}\chronologyannotationreach{John.4.50}{John}{true}\chronologyannotationreach{John.4.51}{John}{true}\chronologyannotationreach{John.4.52}{John}{true}\chronologyannotationreach{John.4.53}{John}{true}\chronologyannotationclaim{composition.gospel-of-john}{composition}{catholic-traditional-v1}{disputed}{96 or one of the succeeding years}{A.D. 96--100}.}%
}
% offertory: Offertory
\expandafter\def\csname triptychchronologyannotation@offertory\endcsname{%
  \chronologyannotationgroup{historical-setting}{preferred}{\textbf{Historical setting}: \chronologyannotationreach{Ps.136.1}{Ps.136}{true}\chronologyannotationclaim{israel.exile.psalm-136-captive-lament}{historical-setting}{catholic-traditional-v1}{preferred}{For the time of Jeremias, and the captivity of Babylon}{The Babylonian captives' lament, after Jerusalem's destruction under Sedecias (derived)}.}%
  \space \chronologyannotationgroup{historical-setting}{preferred}{\textbf{Historical setting}: \chronologyannotationreach{Ps.136.1}{Ps.136}{true}\chronologyannotationclaim{israel.exile.seventy-years}{historical-setting}{catholic-traditional-v1}{preferred}{seventy years}{The seventy years of servitude to Babylon, ``seventy years'' (duration)}.}%
  \space \chronologyannotationgroup{composition}{preferred}{\textbf{Composition}: \chronologyannotationreach{Ps.136.1}{Ps}{true}\chronologyannotationclaim{critical.psalms.latest-composition-boundary}{composition}{catholic-critical-v1}{preferred}{before the Maccabean period, around 165 B.C.; no individual psalm can be dated securely}{Before c. 165 B.C.}}%
  \space \chronologyannotationanchorgroup{israel.exile.psalm-136-captive-lament}{historical-setting}{catholic-traditional-v1}{For the time of Jeremias, and the captivity of Babylon}{israel.exile.third-captivity}{after}{context}{\textbf{Historical background: The third captivity of Juda: the destruction of Jerusalem under Sedecias}: Preferred: \chronologyannotationanchorclaim{israel.exile.third-captivity}{0}{catholic-traditional-v1}{preferred}{passage.george-leo-haydock.douay-rheims-with-haydock-commentary.2014-loreto-feeney-memorial.psalm-70-captivities-usher-chronology}{reported-excluded}{A.M. 3416}{Haydock reporting Ussher; A.M. is not converted}{A.M. 3416 (Haydock reporting Ussher; A.M. is not converted)}; alternatives: \chronologyannotationanchorclaim{israel.exile.third-captivity}{1}{catholic-traditional-v1}{alternate}{artifact.catholic-encyclopedia.volume-3.new-york-1908.newadvent-03315a-9cb437e2,artifact.catholic-encyclopedia.volume-14.new-york-1912.newadvent-14499a-82b12b6e}{traditional-catholic}{586 B.C.}{Reid and Meistermann, Catholic Encyclopedia}{B.C. 586 (Reid and Meistermann, Catholic Encyclopedia)}, \chronologyannotationanchorclaim{israel.exile.third-captivity}{2}{catholic-traditional-v1}{alternate}{artifact.catholic-encyclopedia.volume-8.new-york-1910.newadvent-08652a-189e1a8f}{traditional-catholic}{the fall of Jerusalem (587 B.C.)}{Schets, Catholic Encyclopedia}{B.C. 587 (Schets, Catholic Encyclopedia)}, \chronologyannotationanchorclaim{israel.exile.third-captivity}{3}{catholic-traditional-v1}{alternate}{artifact.catholic-encyclopedia.volume-8.new-york-1910.newadvent-08654a-645bba6c}{reported-traditional}{B.C. 588}{Petavius table reported by Sloet; retained transcription checked, printed-page verification still pending}{B.C. 588 (Petavius table reported by Sloet; retained transcription checked, printed-page verification still pending)}. The represented event is after this historical event. These dates belong to that historical event, not to the passage or its composition. Both events could occur in the same year.}%
}
% communion: Communion
\expandafter\def\csname triptychchronologyannotation@communion\endcsname{%
  \chronologyannotationgroup{traditional-attribution}{disputed}{\textbf{Traditional attribution} -- disputed: \chronologyannotationreach{Ps.118.49}{Ps.118}{true}\chronologyannotationreach{Ps.118.50}{Ps.118}{true}\chronologyannotationclaim{israel.monarchy.david-traditional-era}{traditional-attribution}{catholic-traditional-v1}{disputed}{reigned from 1055 to 1015 B.C.}{David (reign in the usual chronology), B.C. 1055--1015}.}%
  \space \chronologyannotationgroup{historical-setting}{research-pending}{\textbf{Historical setting} -- Occasion date unresolved.}%
  \space \chronologyannotationgroup{composition}{preferred}{\textbf{Composition}: \chronologyannotationreach{Ps.118.49}{Ps}{true}\chronologyannotationreach{Ps.118.50}{Ps}{true}\chronologyannotationclaim{critical.psalms.latest-composition-boundary}{composition}{catholic-critical-v1}{preferred}{before the Maccabean period, around 165 B.C.; no individual psalm can be dated securely}{Before c. 165 B.C.}}%
}


\section{Scriptural Date and Location}

\zlabel{triptych:concise:chronology:start}
\begin{concisedossiertable}
Alleluia & Ps 107:2 (Heb.\ 108) & David, by the title; the psalm repeats the
closes of Pss~56 and 59 &
\chronodate{alleluia}{\chronologyannotation{alleluia}}\\
\cmidrule(lr){1-4}
\dossierprose{Titled \latin{Canticum Psalmi ipsi David}, with no occasion. The
reign is David's in the ``usual chronology'' of the \work{Catholic
Encyclopedia}'s ``David'' (vol.~4, 1908), which reports a later dating beside
it, a reference point and not the psalm's date; the composition date is the
\work{New American Bible, Revised Edition}'s bound for the Psalter.}
\midrule
Communion & Ps 118:49--50 (Heb.\ 119) & The psalmist, ``a sojourner on the
earth'' (v.~19), in the strophe \latin{Zain} &
\chronodate{communion}{\chronologyannotation{communion}}\\
\cmidrule(lr){1-4}
\dossierprose{Titled \latin{Alleluia}; the Davidic attribution is disputed, and the
record holds no occasion. The Introit's verse, Ps~118:1, is dated below.}
\midrule
Offertory & Ps 136:1 (Heb.\ 137) & The captives of Juda by the rivers of
Babylon, after the fall of Jerusalem &
\chronodate{offertory}{\chronologyannotation{offertory}}\\
\cmidrule(lr){1-4}
\dossierprose{Titled \latin{Psalmus David, Ieremiae}. The record gives a setting, not an author: ``For the time of Jeremias, and
the captivity of Babylon'' (Haydock), and the seventy years of servitude (Jer
25:11--12), a duration. The background dates belong to Jerusalem's destruction
under Sedecias, not to the psalm; the year from the creation of the world is
Haydock's report of Ussher, unconverted.}
\midrule
Gradual & Ps 144:15--16 (Heb.\ 145) & David, by the title; God's kingship and
his care for every creature &
\chronodate{gradual}{\chronologyannotation{gradual}}\\
\cmidrule(lr){1-4}
\dossierprose{Titled \latin{Laudatio ipsi David}, without an occasion; the
reign is a reference point, the composition date the Psalter's bound.}
\midrule
Introit & Dan 3:31, 29 and 35, as printed; the verse Ps 118:1 (Heb.\ 119) &
The furnace in the province of Babylon, under Nabuchodonosor (Dan 3:1, 20) &
\chronodate{introit}{\chronologyannotation{introit}}\\
\cmidrule(lr){1-4}
\dossierevent{Narrated event: the three young men, cast into the furnace for
refusing to adore the king's statue, walk in the flame and Azarias prays
(Dan 3:1--25); the narrative gives no year.}
\cmidrule(lr){1-4}
\dossierprose{Daniel and the verse are dated differently, so the cell holds no
single date. The \work{Catholic Encyclopedia}'s ``Book of Daniel'' (vol.~4,
1908) ascribes the Prayer of Azarias, without a year, to an unknown Jewish
author, ``most likely long after the Exile''; St Jerome reports that it is not
in the Hebrew. For the
verse, the record gives the Psalter's composition bound, Ps~118:1: before
c.~165~B.C.; and, as the reference point of the disputed attribution to David,
his reign in the usual chronology, Ps~118:1: 1055--1015~B.C.}
\midrule
Gospel & Jn 4:46--53 & Writing: St John the Apostle, by the Missal's heading.
Event: Cana and Capharnaum, in Galilee (4:46) &
\chronodate{gospel}{\chronologyannotation{gospel}}\\
\cmidrule(lr){1-4}
\dossierevent{Narrated event: ``the second miracle that Jesus did, when he was
come out of Judea into Galilee'' (4:54); the word is spoken at Cana, and the
fever leaves the boy at Capharnaum at the same hour (4:52).}
\cmidrule(lr){1-4}
\dossierprose{The event's year is derived from the reckoning from the founding of Rome
in A.~J. Maas's ``Jesus Christ'' (\work{Catholic Encyclopedia}, vol.~8, 1910);
the two composition dates, both disputed, are from its ``The New Testament''
(vol.~14) and ``The Gospel of Saint John'' (vol.~8).}
\midrule
Epistle & Eph 5:15--21 & Paul, ``a prisoner in the Lord'' (4:1), in Caesarea or
Rome; to ``the saints who are at Ephesus'' (1:1) &
\chronodate{epistle}{\chronologyannotation{epistle}}\\
\cmidrule(lr){1-4}
\dossierprose{Traditional attribution: St Paul, who names himself (1:1) and
writes as ``an ambassador in a chain'' (6:20). The disputed range is the
\work{Catholic Encyclopedia}'s ``Epistle to the Ephesians'' (vol.~5, 1909), for
which Caesarea or Rome ``is still a much mooted question'', and the disputed
year its ``St.\ Paul'' (vol.~11, 1911). Shown last is the modern critical horizon:
the \work{New American Bible, Revised Edition}'s later disciple, around
A.D.~80--100, one possible explanation and not the default answer.}
\end{concisedossiertable}
\zlabel{triptych:concise:chronology:end}
